\section{Conclusion}\label{sec:conclusion}

A retail payment that uses a fiat-backed stablecoin appears cheap by the standard of the announced on-chain transfer fee but carries a structural exposure to two risks that no other retail payment instrument shares: counterparty hazard of the issuer, and a depeg on the secondary market when convertibility at par is interrupted under stress. We read the counterparty component from market proxies and price the stress depeg with a closed-form global game of holder coordination, separately for USDC and for USDT, obtaining holding-leg expected losses of 21 and 60 basis points respectively, both counterparty-led but separated by the severity of the stress depeg, deep for USDC, whose primary redemption is vulnerable to suspension, and mild for USDT, whose primary redemption stayed open. Adding these expected losses and the conversion fees of the circuit to the announced transfer fee and benchmarking against publicly operated instant payment systems, correspondent banking, and fintech remittance providers across corridors in Brazil, the United States, and the European Union, we find that the stablecoin path is dominated where a public instant payment system already serves the corridor, wins where none does, and would be absorbed by an upgraded public cross-border alternative once one is operational. The recent expansion of stablecoin payment use on cross-border corridors is best read as a temporary bridge in payments innovation while public infrastructure is developed, rather than a permanent technological frontier.

The framework applies to any corridor for which the cost data and the calibration episodes are available; the three corridors here span the policy-relevant space of domestic retail, regional remittance, and cross-border flows. Asian payment systems are at the most advanced stage of cross-border interlinkage, and extending the comparison there requires calibration to the corresponding stablecoin pairs, local rate curves, and on-ramp arrangements, which we leave for future work.

Three further extensions are open. A dynamic adoption model would characterise the speed of substitution when a public alternative goes live and the path along which the stablecoin's market share contracts on absorbed corridors. A multi-currency extension would examine whether the substitution effect on the EU-Brazil corridor would change if a euro-denominated stablecoin became liquid at scale. An issuer-side strategic-response extension would endogenise the reserve, asset, and hazard parameters: a contraction in the stablecoin's market following the launch of a public alternative would induce balance sheet reoptimisation, and the structural expected loss itself would respond.
